\documentclass[9pt,lineno]{elife}

\usepackage{amsfonts}

\setcounter{secnumdepth}{3}
\titleformat{\section}
  {\color{eLifeMediumGrey}\Large\bfseries}
  {\thesection}{0.5em}{#1}[]
\titleformat{\subsection}
  {\large\bfseries}
  {\thesubsection}{0.5em}{#1}[]
\titleformat{\subsubsection}
  {\large}
  {\thesubsubsection}{0.5em}{#1}[]
\renewcommand{\arraystretch}{1.3}

\title{Bracket Norm Tracks Causally Important Brain Regions From Population Geometry}

\author[1]{Elliot Tower}
\affil[1]{Independent Researcher}
\corr{elliot@elliottower.ai}{ET}
\contrib{Elliot Tower (\href{https://orcid.org/0000-0001-7004-8884}{0000-0001-7004-8884}): Conceptualization, Methodology, Software, Formal Analysis, Investigation, Data Curation, Writing -- Original Draft, Writing -- Review \& Editing, Visualization.}

\begin{document}
\maketitle

\begin{abstract}
Geometric properties of neural population activity---decoding accuracy, dimensionality, persistent homology, and 16 others---all fail to predict which brain regions are causally necessary for binary choice once the number of recorded neurons is controlled. We demonstrate this confound against optogenetic silencing ground truth (12 regions) and propose \emph{bracket norm}, the evidence-dependent rotation of the choice displacement vector, as the one metric that survives. Raw bracket norm scales as $\sqrt{n}$; the intensive quantity $\text{BN}/\!\sqrt{n}$ eliminates this dependence (CV $2.0\%$ across $25\times$ subsampling). The top-3 regions by $\text{BN}/\!\sqrt{n}$ match optogenetic silencing rankings exactly ($p = 0.012$, exact combinatorial), and the silencing ranges of the BN-defined top and bottom tails do not overlap (permutation $p = 0.0008$, survives Bonferroni). On a held-out dataset (IBL, different task), the correction removes the neuron-count confound and reveals a silencing correlation from non-significant ($\rho = -0.286$, $p = 0.49$) to $\rho = -0.667$ ($p = 0.083$), consistent with task-specific computation. ALM photoinhibition \emph{increases} $\text{BN}/\!\sqrt{n}$ while degrading behavior (32/32 sessions, $p < 10^{-6}$). The neuron-count confound affects all geometric predictors tested; the intensive correction is necessary for valid cross-region comparisons.
\end{abstract}

\noindent\textbf{Impact statement:} Standard geometric measures of neural population activity are confounded by the number of recorded neurons; a dimensionally corrected metric survives and predicts optogenetic silencing effects across brain regions, tasks, and species.

\section{Introduction}

A central goal of systems neuroscience is identifying which brain regions causally drive behavior, and a growing literature uses the \emph{geometry} of population activity to do so \citep{cunningham2014dimensionality, gallego2017neural}. Regions are ranked by decoding accuracy, representational dimensionality, manifold curvature, or equivariance, on the implicit assumption that geometric richness reflects causal involvement. This assumption has not been tested against interventional ground truth.

We test it directly against optogenetic silencing of 12 brain regions during a visual decision task \citep{steinmetz2019distributed, zatkahaas2021}. Brain region abbreviations follow the Allen Common Coordinate Framework (CCF); common equivalents are noted at first use (e.g., ORB = OFC, VISp = V1, ACA = ACC). The basic observable is the \emph{choice displacement vector}---the difference in mean population activity between left-choice and right-choice trials---computed separately at low and high sensory evidence. All 19 geometric metrics we test predict silencing importance only because they track the number of recorded neurons; after controlling for recording yield, all 19 collapse. We propose \emph{bracket norm}, which measures how much the choice displacement vector rotates as evidence changes, and show it survives the confound. Because raw bracket norm inherits a $\sqrt{n}$ dimensional scaling, we derive the corrected intensive quantity $\text{BN}/\!\sqrt{n}$, which eliminates the neuron-count dependence. On a held-out dataset (IBL wheel-turning), the correction eliminates the neuron-count confound and uncovers a silencing prediction of $\rho = -0.667$ ($p = 0.083$), masked in the raw data ($\rho = -0.286$, $p = 0.49$), with region rankings that are task-specific: licking-task and wheel-turning-task rankings are uncorrelated ($\rho = -0.079$, $n = 18$), and the sign of the silencing correlation flips---an observation we interpret post-hoc as task specificity (\S\ref{sec:scope}).

Directly perturbing a region (ALM photoinhibition) \emph{increases} $\text{BN}/\!\sqrt{n}$ while degrading behavior, suggesting the metric tracks per-neuron geometric expressiveness rather than computational quality. On human ECoG during speech production, the metric ranks articulatory motor cortex highest with zero electrode-count confound (\S\ref{sec:ecog}).

%% ========== RESULTS ==========

\section{Results}

\subsection{All 19 geometric metrics are confounded by neuron count}
\label{sec:confound}

We computed 19 geometric metrics for each of the 12 silencing regions---established measures from the neural geometry literature (decoding accuracy, dimensionality, participation ratio, persistent homology, curvature, equivariance, Sliced Wasserstein distance) and additional geometric quantities we constructed (cloud variance, force alignment, commutativity, transport metrics)--- and correlated each with silencing effect size (the behavioral deficit on inactivation). For each metric we report raw Spearman $\rho$, its correlation with neuron count, and the partial $\rho$ controlling for neuron count.

Recording yield predicts silencing at $\rho = +0.762$ and correlates with nearly every geometric metric at $\rho > 0.8$. After partialing neuron count, all 19 collapse (Table~\ref{tab:confound}). We then propose bracket norm, a new metric not previously used in this literature, which retains substantial predictive power. In joint OLS regression, bracket norm alone explains more variance in silencing importance ($R^2 = 0.652$) than neuron count alone ($R^2 = 0.580$); the two carry partially overlapping but not redundant information (full model $R^2 = 0.809$; Appendix~\ref{app:joint}).

\begin{table}[t]
\centering
\caption{Geometric metrics versus optogenetic silencing importance. All 19 collapse after controlling for neuron count; bracket norm (proposed here, top row) survives. Metrics include both established measures from the literature and geometric quantities constructed for this comparison. Thirteen of 19 shown; 6 additional variants omitted for space; all follow the same pattern.}
\label{tab:confound}
\begin{tabular}{lccc}
\toprule
Metric & Raw $\rho$ & Corr.\ with neuron count & Partial $\rho$ \\
\midrule
\textbf{Bracket norm (decision window)} & \textbf{+0.883} & \textbf{+0.850} & \textbf{+0.753} \\
Decision-window SEIS & +0.706 & +0.902 & +0.067 \\
Persistent homology halflife & +0.650 & +0.825 & +0.058 \\
LDA accuracy & +0.650 & --- & --- \\
ER (normalized) & $-0.650$ & $-0.909$ & +0.158 \\
Curvature index & +0.615 & +0.842 & $-0.112$ \\
SEIS equivariance (full trial) & +0.601 & +0.867 & $-0.185$ \\
Cloud variance & +0.573 & +0.881 & $-0.321$ \\
Force alignment & +0.545 & --- & +0.015 \\
Transport metrics & +0.531 & --- & $-0.376$ \\
Commutativity (full trial) & +0.524 & --- & +0.071 \\
Noise variance & +0.510 & --- & $-0.414$ \\
Sliced Wasserstein & +0.140 & --- & --- \\
Entropy production & +0.042 & --- & --- \\
\bottomrule
\end{tabular}
\end{table}

Splitting the 12 regions at median effective rank (a neuron-count proxy), the headline SEIS correlation ($\rho = +0.706$) is entirely \emph{between} strata. Within the high-alpha stratum $\rho = -0.029$; within the low-alpha stratum $\rho = +0.086$. SEIS equivariance appears to add no predictive power once dimensionality is held fixed.

Neuron count may dominate because it is the least noisy measurement, and any variable with large dynamic range will out-correlate noisy geometric metrics at $n = 12$. The sign-flip table rules this out (Table~\ref{tab:signflip}): after partialing neuron count, several metrics do not merely collapse to zero---they \emph{reverse sign coherently}. Less noise, less persistence, and less transport spread become associated with \emph{more} causal importance. Collapse-to-zero would indicate noise dominance; directional reversal suggests real structure. The sign pattern is consistent with bracket norm: regions doing active computation may have compact, evidence-sensitive choice representations rather than diffuse, redundant ones.

\begin{table}[t]
\centering
\caption{Metrics that reverse sign after controlling for neuron count. Coherent reversal rules out noise dominance as an explanation for the confound.}
\label{tab:signflip}
\begin{tabular}{lcc}
\toprule
Metric & Raw $\rho$ & Partial $\rho$ \\
\midrule
$H_0$ persistence & $-0.053$ & $-0.532$ \\
Bracket norm (full trial) & $+0.343$ & $-0.460$ \\
Noise variance & $+0.510$ & $-0.414$ \\
Transport metrics & $+0.531$ & $-0.376$ \\
Betti-0 at 25\% & $+0.007$ & $+0.364$ \\
\bottomrule
\end{tabular}
\end{table}

Neuron count may instead be a mediator---causally important regions \emph{recruit more neurons}---in which case partialing it destroys real signal. The synthetic subsampling test (\S\ref{sec:dim_correction}) rules this out: subsampling from the \emph{same recording} at $n \in \{10, \ldots, 250\}$ produces $\rho = +0.989$ between BN and $n$, despite identical underlying biology. The dependence is a recording-side artifact, not a biological recruitment effect. We argue that any geometric predictor of causal importance in Neuropixels data inherits this confound unless explicitly controlled, rendering raw correlations difficult to interpret.

\subsection{Bracket norm tracks silencing importance}
\label{sec:bracket}

Bracket norm measures the evidence-dependent rotation of the choice displacement vector during the decision window (\S\ref{sec:bn_def}; Figure~\ref{fig:quiver}).

\begin{figure}[t]
\centering
\includegraphics[width=\textwidth]{figures/fig_quiver.pdf}
\caption{Kernel-smoothed choice displacement field for PL (prelimbic cortex, left) and VISam (anteromedial visual area, right), computed from Steinmetz et al.\ (2019) spike counts during a late decision subwindow (250--350\,ms post-stimulus; the full decision window used for all analyses is 150--350\,ms). Each arrow shows the local choice displacement vector at a grid point in (sensory evidence $\times$ choice-axis projection) space; color encodes per-neuron gain ($\lVert\xi\rVert / \sqrt{n}$). In PL, arrows are large and colored (high gain), with orientation varying across the evidence axis---evidence reshapes the choice code. In VISam, arrows are uniformly small and dark---evidence does not modulate the choice signal. PL ranks \#1 by silencing effect and \#2 by bracket norm (BN\,=\,0.323); VISam ranks near the bottom (BN\,=\,0.065), a $5\times$ contrast.}
\label{fig:quiver}
\end{figure}

Bracket norm predicts silencing at raw $\rho = +0.883$ ($p = 0.003$, exact permutation test) and partial $\rho = +0.753$ after controlling for neuron count ($n = 9$, the regions with sufficient trials per evidence quartile). After $\sqrt{n}$ normalization (\S\ref{sec:dim_correction}), the rank correlation persists with the neuron-count confound removed: $\text{BN}/\!\sqrt{n}$ predicts silencing at $\rho = +0.833$ ($p = 0.007$; bootstrap 95\% CI: $[0.26,\; 1.00]$, $10{,}000$ resamples; Figure~\ref{fig:ranking}).

Applying Bonferroni correction for 20 simultaneous tests (19 comparison metrics plus bracket norm), the corrected threshold is $\alpha = 0.0025$. The non-overlapping tail test (permutation $p = 0.0008$) survives this correction. The rank correlations ($p = 0.003$ raw, $p = 0.007$ normalized) and partial $\rho$ ($p = 0.031$) do not survive Bonferroni individually, but they address complementary questions---whether bracket norm predicts at all, whether it predicts beyond neuron count---and the combinatorial test, which makes no distributional assumptions, confirms the result at the corrected threshold.

The top-3 bracket-norm regions (ORB, PL, VISpm) are identical to the top-3 by silencing ($p = 0.012$, exact combinatorial), and the silencing ranges of the BN-defined top and bottom tails do not overlap ($\min(\text{top-3}) = 0.225 > \max(\text{bottom-3}) = 0.145$; permutation $p = 0.0008$, survives Bonferroni; Table~\ref{tab:tails}, Figure~\ref{fig:ranking}). The top regions---orbitofrontal cortex (ORB/OFC), prelimbic cortex (PL), and posteromedial visual area (VISpm)---are association areas involved in evidence evaluation. The bottom regions by bracket norm---anteromedial visual area (VISam), anterior cingulate cortex (ACA/ACC), and primary visual cortex (VISp/V1)---are early sensory or monitoring areas with low silencing effects. The result is robust: $\rho$ is positive across all 8 temporal windows tested, all 9 jackknife estimates are positive, and pre-stimulus bracket norm already predicts silencing ($\rho = +0.533$; Appendix~\ref{app:robustness}). The signal is stable across sessions: the mean per-region bracket norm (averaged across sessions) predicts silencing at $\rho = +0.618$ ($p = 0.043$, $n = 11$ regions with $\geq 2$ sessions), while per-region variability does not ($\rho = +0.045$, $p = 0.89$).

\begin{table}[t]
\centering
\caption{Region rankings by bracket norm versus optogenetic silencing. Top-3 sets are identical ($p = 0.012$, exact combinatorial); silencing ranges of top and bottom tails do not overlap (permutation $p = 0.0008$).}
\label{tab:tails}
\begin{tabular}{clclc}
\toprule
Rank & By bracket norm & (BN) & By silencing & (effect) \\
\midrule
1 & ORB & 0.188 & PL & 0.333 \\
2 & PL & 0.172 & ORB & 0.309 \\
3 & VISpm & 0.163 & VISpm & 0.225 \\
7 & VISam & 0.092 & RSP & 0.142 \\
8 & ACA & 0.073 & VISp & 0.141 \\
9 & VISp & 0.067 & VISam & 0.082 \\
\bottomrule
\end{tabular}
\end{table}

\begin{figure}[t]
\centering
\includegraphics[width=0.5\textwidth]{figures/fig_ranking_scatter.pdf}
\caption{$\text{BN}/\!\sqrt{n}$ predicts optogenetic silencing effect across 9 brain regions ($\rho = +0.833$, $p = 0.007$, exact permutation test). Red: top-3 regions by both metrics (ORB, PL, VISpm); blue: bottom-3 by bracket norm (VISam, ACA, VISp). The green band marks the separation gap---silencing ranges of the two groups do not overlap (permutation $p = 0.0008$). Dashed line: OLS fit.}
\label{fig:ranking}
\end{figure}

\subsection{Photoinhibition increases bracket norm while degrading behavior}
\label{sec:mechanism}

As an independent test of whether bracket norm tracks causal importance, we examined a region with direct interventional data: ALM photoinhibition (32 sessions; \citealp{svoboda_alm}). Photoinhibition \emph{increases} bracket norm by 68\% (control $\text{BN}/\!\sqrt{n} = 0.34 \to$ photostim $0.57$, 32/32 sessions, Wilcoxon $p < 10^{-6}$; Figure~\ref{fig:photostim}) while reducing accuracy from 86.8\% to 74.7\%. If bracket norm measured ``amount of computation,'' disruption should lower it; instead the disrupted circuit has \emph{more} geometric structure. An independent replication on ALM photoinhibition data \citep{dandi000007, dandi000009} (pole-position discrimination) confirms the direction: photostimulation increases $\text{BN}/\!\sqrt{n}$ by 59\% (control $0.199 \pm 0.087$, $n = 52$; stim $0.317 \pm 0.098$, $n = 46$; Cohen's $d = +1.27$). The rotation angle between high- and low-evidence choice vectors barely changes ($1.60 \to 1.55$\,rad, $-3\%$), suggesting a dissociation between coding \emph{gain} and coding \emph{direction}: disruption appears to amplify per-neuron response without reorienting the coding axis, consistent with reduced recurrent suppression exposing stronger feedforward drive. In the replication, the principal angle between control and photostim bracket vectors is $85.1^\circ \pm 17.7^\circ$ (median $84.5^\circ$, $n = 44$), while within-condition rotation is preserved (ctrl $81.0^\circ$, stim $79.9^\circ$)---perturbation changes \emph{which neurons} carry the signal without disrupting \emph{how} evidence modulates choice encoding.

\begin{figure}[t]
\centering
\includegraphics[width=0.45\textwidth]{figures/fig_photostim_slopegraph.pdf}
\caption{ALM photoinhibition increases $\text{BN}/\!\sqrt{n}$ in every session. Each line connects one session's control and photostimulation values ($n = 32$ paired sessions, Wilcoxon signed-rank $p < 10^{-6}$, mean $+68\%$: control $0.34 \pm 0.01$, photostim $0.57 \pm 0.02$). Data from \citet{svoboda_alm}; DANDI:000007, 000009 \citep{dandi000007, dandi000009}.}
\label{fig:photostim}
\end{figure}

We tested three alternative explanations. Z-score artifact: rejected---raw unnormalized bracket norm also rises ($10.6 \to 17.6$, 30/32 sessions). Subsampling: rejected---random dropout of matched neuron counts \emph{decreases} bracket norm ($2.10 \to 1.50$ at 50\%), and targeted dropout of the specific suppressed neurons produces $\text{BN} = 2.16$ (41st percentile of the random distribution). The regime-shift interpretation survives every test we applied (Table~\ref{tab:mechanism}).

\begin{table}[t]
\centering
\caption{Twelve pre-specified tests of the recurrent-gain-reduction mechanism. ALM photoinhibition shifts the network from recurrence-dominated to input-dominated dynamics. The autocorrelation timescale $\tau$ increases, dissociating slow recurrent dynamics from the fast evidence-dependent gain bracket norm captures.}
\label{tab:mechanism}
\small
\begin{tabular}{llc}
\toprule
Test & Result (control $\to$ photostim) & Status \\
\midrule
Targeted dropout reproduces effect & targeted 2.16, random 2.18, stim 3.81 & rejected (not subtraction) \\
Participation ratio & $23.6 \to 18.0$, 0/32 increase, $p<10^{-6}$ & confirmed \\
Per-neuron selectivity (auROC) & $0.523 \to 0.540$, 32/32 sessions & confirmed \\
Across-quartile separation & $0.047 \to 0.145$ (triples) & confirmed \\
Within-quartile variance & $0.867 \to 0.775$ & confirmed \\
Eigenspectrum slope & $-1.31 \to -2.10$, 33/48 more low-rank & confirmed \\
Top-3 eigenvalue fraction & $0.547 \to 0.583$ & confirmed \\
Fano factor (single-neuron noise) & $7.83 \to 7.87$ (unchanged) & confirmed \\
Temporal BN profile & $1.7\times$ at every timepoint incl.\ pre-stim & confirmed \\
BN peak latency & $1350 \to 1100$\,ms ($-250$\,ms) & confirmed \\
BN increase vs.\ accuracy decrease & $\rho = -0.059$, $p = 0.75$ (independent) & confirmed \\
Autocorrelation timescale $\tau$ & $6.26 \to 6.40$\,ms ($p = 0.001$, opposite direction) & informative \\
\bottomrule
\end{tabular}
\end{table}

The autocorrelation timescale $\tau$ responds in the opposite direction from bracket norm (Table~\ref{tab:mechanism}), and the two quantities are orthogonal across regions ($\rho = +0.231$, $p = 0.17$). Simulated neuron loss supports this interpretation: randomly dropping neurons \emph{decreases} BN ($2.10 \to 1.50$ at 50\% dropout), and the photoinhibition effect ($\text{BN} = 3.81$) is $5.7\times$ larger than what neuron removal alone produces.

\subsection{Human ECoG generalizes across species and recording modality}
\label{sec:ecog}

Bracket norm was developed and validated entirely on mouse Neuropixels data. To test whether the metric generalizes beyond mice and beyond single-unit recordings, we computed $\text{BN}/\!\sqrt{n}$ on human ECoG recordings during speech production \citep{dandi000019} (19 sessions across 5 participants, 5 cortical areas with $\geq 2$ sessions each). Evidence is operationalized as consonant identity (the articulatory target); choice is the produced syllable.

Ventral sensorimotor cortex (vSMC; $\text{BN}/\!\sqrt{n} = 0.382$) ranks highest, followed by caudal middle frontal ($0.323$), superior temporal gyrus (STG; $0.320$), pars opercularis ($0.315$), and supramarginal gyrus ($0.267$). vSMC, the primary articulatory cortex and a candidate functional homolog of mouse ALM, carries the strongest evidence-dependent choice signal---$1.4\times$ higher than pars opercularis and $1.4\times$ higher than supramarginal gyrus. Classical neurology places Broca's area (pars opercularis) at the center of speech production, but modern functional mapping separates phonological planning (Broca's) from articulatory execution (vSMC). If bracket norm tracks execution-level dynamics, it should rank pars opercularis below vSMC, and it does---a testable prediction that distinguishes the metric from crude ``does this area do speech'' measures.

Raw BN correlates with electrode count at $\rho = +0.800$ across ECoG areas. After normalization, $\text{BN}/\!\sqrt{n}$ shows $\rho = 0.0$ ($p = 1.0$; exact value---for $n = 5$, Spearman $\rho$ takes discrete values and $\Sigma d^2 = 20$ yields $\rho = 0$ exactly) with electrode count. The confound present in raw BN on Neuropixels data also affects raw BN on ECoG, and the same normalization resolves it (Table~\ref{tab:ecog}). In a production task the mapping from consonant identity to syllable is largely deterministic, which may contribute to the near-uniform rotation angles observed across areas ($1.19$--$1.49$\,rad); future work should test the metric in ECoG tasks with explicit sensory uncertainty.

\begin{table}[t]
\centering
\caption{Human ECoG speech production: area rankings by $\text{BN}/\!\sqrt{n}$ and confound test against electrode count. The normalization eliminates the confound present in raw BN.}
\label{tab:ecog}
\small
\begin{tabular}{lcccc}
\toprule
Area & BN (raw) & $\text{BN}/\!\sqrt{n}$ & Rotation angle (rad) & Electrodes \\
\midrule
vSMC & 3.51 & 0.382 & 1.472 & 84 \\
caudal middle frontal & 1.44 & 0.323 & 1.454 & 20 \\
STG & 3.11 & 0.320 & 1.431 & 98 \\
pars opercularis & 1.60 & 0.315 & 1.487 & 26 \\
supramarginal & 1.33 & 0.267 & 1.194 & 25 \\
\midrule
\multicolumn{5}{l}{Confound: $\rho$ vs.\ electrode count} \\
\midrule
 & $+0.800$ & $0.0$ & $+0.100$ & --- \\
\bottomrule
\end{tabular}
\end{table}

\subsection{Region rankings are task-specific}
\label{sec:scope}

Bracket norm was developed on the Steinmetz visual-discrimination task (licking response), where $\text{BN}/\!\sqrt{n}$ predicts silencing importance at $\rho = +0.833$ (Figure~\ref{fig:ranking}; \S\ref{sec:bracket} reports raw BN at $\rho = +0.883$ before $\sqrt{n}$ normalization). On the IBL wheel-turning task (86 sessions, 8 overlapping regions), $\text{BN}/\!\sqrt{n}$ correlates with the \emph{same} Steinmetz silencing ground truth at $\rho = -0.667$ ($p = 0.083$, exact permutation test). The sign reversal---positive within-task, negative cross-task---raises the question of whether the metric is measuring something consistent.

The resolution is task specificity. Across the 18 regions recorded in both datasets, $\text{BN}/\!\sqrt{n}$ rankings are uncorrelated ($\rho = -0.079$, $p = 0.75$, $n = 18$): regions that rank high on the licking task do not rank high on the wheel-turning task. The silencing ground truth measures behavioral deficit during licking. Regions that compute for the licking task score high on the Steinmetz $\text{BN}/\!\sqrt{n}$ ranking and are missed when silenced during licking (positive $\rho$). Regions that compute for the wheel-turning task score high on the IBL $\text{BN}/\!\sqrt{n}$ ranking and are dispensable for licking (negative $\rho$). This task-specificity interpretation is post-hoc---constructed after observing the sign flip---but it generates a falsifiable prediction: optogenetic silencing during \emph{wheel-turning} should produce a positive $\rho$ with IBL $\text{BN}/\!\sqrt{n}$, since the metric would then be evaluated against the task it measures. No such dataset currently exists.

Two pieces of evidence support this interpretation. First, the top IBL regions include sensorimotor junction structures (CL, ZI, APN) and a higher visual area (VISrl, rostrolateral visual cortex)---consistent with a wheel-turning task requiring rapid visual-to-motor conversion---while the top Steinmetz regions (ORB, PL, VISpm) are prefrontal and association areas involved in deliberative evaluation---consistent with a licking task where evidence must be integrated before report (Figure~\ref{fig:cross_task}). Second, the $\sqrt{n}$ correction resolves the IBL neuron-count confound, confirming the geometric signal was present but masked by recording yield. The result is robust to the choice of normalization method (Table~\ref{tab:method_robustness}). We report per-session normalization (Method~B) throughout for consistency with the Steinmetz analysis, where both methods are equivalent.

\begin{table}[t]
\centering
\caption{IBL silencing prediction under two normalization methods ($n = 8$ regions). Method~B (per-session normalization, then average) is the unit-of-observation approach; Method~A (aggregate raw values, then normalize) gives a stronger result due to Jensen's inequality on the concave $1/\!\sqrt{n}$. Both eliminate the neuron-count confound.}
\label{tab:method_robustness}
\small
\begin{tabular}{lcccc}
\toprule
Method & $\rho$ vs.\ silencing & $p$ & $\rho$ vs.\ neuron count \\
\midrule
Raw BN (no correction) & $-0.286$ & $0.49$ & $+0.690$ \\
B: $\text{mean}(\text{BN}/\!\sqrt{n} \text{ per session})$ & $-0.667$ & $0.083$ & $-0.119$ \\
A: $\overline{\text{BN}}/\!\sqrt{\bar{n}}$ & $-0.762$ & $0.028$ & $-0.119$ \\
\bottomrule
\end{tabular}
\end{table}

\begin{figure}[t]
\centering
\includegraphics[width=0.5\textwidth]{figures/fig_cross_task_scatter.pdf}
\caption{$\text{BN}/\!\sqrt{n}$ rankings are task-specific. Each point is a brain region recorded in both the Steinmetz licking task and the IBL wheel-turning task ($n = 18$). Rankings are uncorrelated ($\rho = -0.079$, $p = 0.75$): regions important for one task are not systematically important for the other. Red squares: regions with optogenetic silencing ground truth (silencing measured during the licking task). The sign flip in silencing prediction ($+0.833$ within-task, $-0.667$ cross-task) follows from this task specificity.}
\label{fig:cross_task}
\end{figure}

Leave-one-out cross-validation on the IBL data produces negative correlations for all 8 held-out regions (mean $\rho = -0.656 \pm 0.096$), so no single region drives the result. Six pre-specified diagnostics confirm the underlying geometry is intact: choice vectors rotate with evidence at $0.98$--$1.41$\,rad, and evidence-conditioned choice subspaces are distinct ($34$--$52^\circ$ principal angles; Table~\ref{tab:ibl}). Rotation angle predicts silencing weakly ($\rho = -0.500$, $p = 0.21$), consistent with the Svoboda finding that causal disruption alters coding gain rather than coding direction.

\begin{table}[t]
\centering
\caption{Six pre-specified IBL diagnostics. The underlying geometry is intact; the dynamic range is compressed.}
\label{tab:ibl}
\begin{tabular}{llc}
\toprule
Diagnostic & Result & Verdict \\
\midrule
H32 BN range compression & CV $0.15$ vs.\ Steinmetz $0.45$ ($3\times$ narrower) & compressed \\
H33 Evidence-quartile separation & weak & weak \\
H34 Choice vector rotation & $0.98$--$1.41$\,rad (vectors do rotate) & geometry intact \\
H35 Grassmannian subspace structure & $34.5$--$51.5^\circ$ principal angles & distinct subspaces \\
H36 Per-session BN variance & CV $0.26$--$0.68$ (moderate) & not washing out signal \\
H37 Selectivity vs.\ BN (206 regions) & $\rho = 0.197$, $p = 0.005$ & weak, not explanatory \\
\bottomrule
\end{tabular}
\end{table}

On the Allen Visual Behavior Neuropixels change-detection task (20 sessions, 49 regions, go/no-go licking), two evidence operationalizations produce strongly correlated $\text{BN}/\!\sqrt{n}$ rankings ($\rho = +0.803$, top-10 overlap 7/10), with partial $\rho = +0.757$ controlling neuron count (Table~\ref{tab:allen}). No silencing data exists for these regions. Subcortical sensorimotor junctions (SNr, APN, MRN) rank highest, consistent with a go/no-go task requiring rapid decision-to-action conversion---a different top set than either the Steinmetz ranking (prefrontal) or the IBL ranking (thalamic-hypothalamic), reinforcing the task-specificity of the metric.

\begin{table}[t]
\centering
\caption{Allen VBN: evidence-operationalization robustness (49 regions, 20 sessions).}
\label{tab:allen}
\small
\begin{tabular}{lc}
\toprule
Evidence operationalization & $\rho$ \\
\midrule
Binary (is\_change) vs.\ neuron count, 7/20 sessions, 37 regions & $+0.767$ \\
Trial-length (continuous) vs.\ neuron count, 20/20 sessions, 49 regions & $+0.703$ \\
Per-image detection rate vs.\ neuron count, 18/20 sessions, 47 regions & $+0.530$ \\
\midrule
Trial-length vs.\ per-image ranking, partial $\rho$ (ctrl NC) & $\mathbf{+0.757}$ \\
\bottomrule
\end{tabular}
\end{table}

We tested whether coarse anatomical category predicts $\text{BN}/\!\sqrt{n}$ ranking using a pre-registered analysis (hypotheses, category assignments, and statistical tests frozen at commit \texttt{7719b64} before any categorical tests were performed; Table~\ref{tab:anatomical}). Across the full IBL dataset (43 regions with $\geq 5$ observations, 6 Allen CCF divisions), the predicted category ordering does not hold, category membership explains near-zero variance, and fiber-tract regions score comparably to neural tissue. The category ordering is unstable across sessions, holding in only 1 of 75 leave-one-session-out folds. All results are robust across inclusion thresholds ($\geq 5$, $\geq 8$, $\geq 10$ observations).

\begin{table}[t]
\centering
\caption{Pre-registered anatomical category hypotheses (H1--H6). All hypotheses, category assignments, and pass thresholds frozen at commit \texttt{7719b64} before analysis.}
\label{tab:anatomical}
\small
\begin{tabular}{llcc}
\toprule
ID & Hypothesis & Result & Verdict \\
\midrule
H1 & Category ordering (midbrain $>$ isocortex $>$ hippo.\ $>$ thalamus) & JT $p = 0.34$ & fail \\
H2 & Fiber tracts score below neural tissue (negative control) & MW $p = 0.32$ & fail \\
H3 & Sensorimotor nuclei enriched in top-5 (CL, ZI, APN) & hyper.\ $p = 0.024$, BH $= 0.048$ & pass \\
H4 & Isocortex gradient (motor $>$ association $>$ sensory) & exploratory & --- \\
H5 & Category explains variance ($\eta^2$) & $\eta^2 \approx 0$, KW $p = 0.77$ & fail \\
H6 & Category ordering stable across sessions (LOO) & 1/75 folds (1.3\%) & fail \\
\bottomrule
\end{tabular}
\end{table}

The category-level null results are informative. The predicted ordering (midbrain $>$ isocortex $>$ hippocampal $>$ thalamus) fails because thalamus has the \emph{highest} category median ($0.374$), the opposite of the predicted lowest (Jonckheere-Terpstra $p = 0.34$, $10^5$ permutations). That fiber-tract regions---axon bundles with no local cell bodies---score comparably to neural tissue (Mann-Whitney $p = 0.32$) likely reflects the fact that Neuropixels electrodes traversing fiber tracts still record action potentials from axons of passage, which inherit the evidence-dependent coding properties of their somatic populations.

Individual sensorimotor transformation nuclei do concentrate at the top of the ranking: 3 of the top-5 regions (CL, ZI, APN) are sensorimotor junction structures spanning thalamus, hypothalamus, and midbrain respectively (hypergeometric $p = 0.024$, Benjamini-Hochberg corrected $p = 0.048$). These regions share a functional role---converting graded sensory evidence into motor commands for wheel-turning---despite belonging to three different anatomical divisions. The dissociation between category-level nulls and region-level enrichment indicates that $\text{BN}/\!\sqrt{n}$ tracks functional position in the evidence-to-action pathway rather than anatomical address.

%% ========== DISCUSSION ==========

\section{Discussion}

Evidence-dependent rotation of choice encoding predicts which brain regions are causally important for binary choice, succeeding where all 19 geometric metrics tested here---both established and constructed---fail to predict independently of recording yield. The dimensionally corrected quantity $\text{BN}/\!\sqrt{n}$ eliminates the residual neuron-count scaling, removing the neuron-count confound on the IBL data and revealing a silencing prediction ($\rho = -0.667$) masked in the raw data ($\rho = -0.286$), and enabling confound-free generalization to human ECoG. Region rankings are task-specific: licking-task and wheel-turning-task rankings are uncorrelated ($\rho = -0.079$, $n = 18$), and each task's ranking predicts silencing importance for the appropriate motor output. The metric captures which regions are computing for a given task, a property that varies with behavioral demands.

We interpret the photoinhibition results through a recurrent attractor model of ALM. Photoinhibition reduces excitatory drive, plausibly shifting the network from recurrence-dominated to input-dominated dynamics: external evidence gains more leverage over a lower-dimensional state. This account predicts the observed pattern---higher bracket norm, lower dimensionality, higher per-neuron selectivity, and lower accuracy (integration requires recurrence). The $1.7\times$ elevation at \emph{every} timepoint, including 150\,ms before stimulus, suggests a tonic change in operating regime rather than a phasic processing effect. Under this interpretation, expressive regions are causally important because the animal depends on their structured representation; within a region, disrupting recurrence amplifies expressiveness at the cost of function. The gain/direction dissociation---disrupting recurrence raises $\text{BN}/\!\sqrt{n}$ without rotating the coding axis---suggests that bracket norm tracks a specific computational quantity, the leverage of evidence on choice representation, rather than a generic measure of neural activity. A radial--angular decomposition of the bracket vector confirms that cross-region variation in bracket norm is primarily a gain (magnitude) effect rather than a rotation (curvature) effect (Appendix~\ref{app:decomposition}).

We also measured autocorrelation timescale $\tau$ as an orthogonal probe of the recurrent-gain account. Photoinhibition slightly increased $\tau$ ($6.26 \to 6.40$\,ms, Wilcoxon $p = 0.001$, $n = 48$ sessions), orthogonal to bracket norm across regions ($\rho = +0.231$, n.s.), indicating that $\tau$ and $\text{BN}/\!\sqrt{n}$ track dissociable timescales of recurrent dynamics. Bracket norm captures fast evidence-dependent gain; $\tau$ reflects slow recurrent modes that photoinhibition does not abolish.

A growing literature ranks brain regions by the geometry of their population activity---decoding accuracy \citep{steinmetz2019distributed}, dimensionality \citep{cunningham2014dimensionality, stringer2019high}, manifold curvature \citep{gardner2022toroidal}, and equivariance \citep{ostrow2023beyond}---on the implicit assumption that geometric richness reflects computational importance. We test this assumption directly against interventional ground truth and find that all 19 metrics we evaluate---both established measures and geometric quantities we constructed---are confounded with the number of recorded neurons. To our knowledge, no prior work tests a comprehensive set of geometric metrics against optogenetic silencing data and reports that nearly all collapse after controlling for recording yield.

Finite-sampling effects on similarity measures are well characterized: CKA sensitivity depends on participation ratio \citep{harvey2024representational}, neuron subsampling underestimates CKA and CCA \citep{murphy2024biased}, and dimensionality estimates carry systematic finite-sample bias \citep{chung2026dimensionality, spectral2025neurips}. Our confound is distinct. These prior results concern estimation error at fixed population size; ours concerns recording yield---regions with more neurons produce higher-dimensional embeddings and proportionally larger Euclidean norms, inflating cross-region importance rankings independently of biological signal. The $\sqrt{n}$ correction we derive addresses this dimensional scaling, not finite-sample noise.

Standard causal methods for time-series data include convergent cross-mapping \citep{sugihara2012detecting}, Granger causality \citep{granger1969investigating}, and conditional-independence methods such as PCMCI \citep{runge2019detecting}. All require long time series (hundreds of points) for Takens embedding or lag estimation. Tangent-space causal inference \citep{butler2024tsci} recovers causal structure from Jacobian predictability without lag estimation, though it still benefits from dense temporal sampling. Trial-structured neuroscience experiments produce 6--10 evidence bins per session---too few for any of these methods (Appendix~\ref{app:causal_baselines}). Bracket norm extracts geometric structure from trial-averaged comparisons natural to cognitive neuroscience experimental designs.

Differential geometry has been applied to neural coding through information geometry \citep{amari2016information} and null-space decompositions \citep{kaufman2014cortical}. The Lie bracket $[\xi_{\text{choice}}, \xi_{\text{evidence}}]$ that bracket norm approximates measures the noncommutativity of two vector fields---how much traversing choice then evidence differs from evidence then choice. Our discrete connection and holonomy analysis (Appendix~\ref{app:holonomy}) makes this interpretation concrete on a choice bundle over evidence space, showing uniform curvature across regions and localizing cross-region variation to gain.

The principal limitation is the size of the interventional ground truth. Region-resolved optogenetic silencing during a quantified perceptual decision task is, to our knowledge, available for 12 regions \citep{steinmetz2019distributed, zatkahaas2021}, of which 9 have sufficient trials per evidence quartile to estimate bracket norm. This sample is not a design choice but a property of the current literature, and it limits the precision of the headline correlation: the bootstrap 95\% confidence interval on $\rho = +0.833$ is wide ($[0.26, 1.00]$; the upper bound reaches $1.00$ because tied ranks in small-$n$ resamples can produce perfect correlation). We therefore base the central claim on rank-based structural tests that remain valid at small $n$---a double dissociation with non-overlapping ranges (permutation $p = 0.0008$) and tail monotonicity with zero counterexamples---rather than on the point estimate alone. As additional silencing datasets become available, the correlation should be re-estimated.

Several further limitations bound the scope of the conclusions. First, raw bracket norm retains a substantial correlation with neuron count ($\rho = +0.850$); $\text{BN}/\!\sqrt{n}$ eliminates this ($\rho = -0.119$ on IBL silencing regions, $\rho = 0.000$ on ECoG), but the normalization assumes z-scored neurons contribute independently---structured noise correlations could violate this. Second, the within-region mechanism is established in a single recurrent region (ALM), and whether the recurrent-gain account generalizes to other regions remains untested. Third, the human ECoG result lacks interventional ground truth; the area ranking is anatomically plausible but not causally validated. Silencing ground truth exists only for the Steinmetz licking task; a second silencing dataset for wheel-turning would directly test the task-specificity claim that IBL $\text{BN}/\!\sqrt{n}$ predicts wheel-task importance positively.

More broadly, geometric predictors of causal importance in neural population data inherit a $\sqrt{n}$ dimensional scaling that should be explicitly corrected. All 19 metrics tested here fail that requirement.

\paragraph{Data availability.} All datasets analyzed are publicly available through the DANDI Archive: Steinmetz Neuropixels recordings (DANDI:000017; \citealp{steinmetz2019distributed, zatkahaas2021}), ALM photoinhibition data (DANDI:000007, 000009; \citealp{svoboda_alm, dandi000007, dandi000009}), and human ECoG speech production (DANDI:000019; \citealp{dandi000019}). The Allen Institute Visual Behavior Neuropixels dataset \citep{siegle2021survey} and the International Brain Laboratory Brain-Wide Map \citep{ibl2025brainwide} are available through their respective data portals. Analysis code, pre-specified hypotheses and pass thresholds, and per-experiment results are available at \url{https://github.com/elliottower/bracket-norm}.

\paragraph{Competing interests.} The author declares no competing interests.

\paragraph{Funding.} No external funding was received for this work.

\paragraph{Ethics.} This study is a secondary analysis of previously published, publicly available datasets. All original studies obtained appropriate ethical approvals: the Steinmetz et al.\ experiments were approved by the UCL Animal Welfare and Ethical Review Body; the Svoboda lab photoinhibition experiments were conducted under IACUC protocols at the Janelia Research Campus; the human ECoG recordings were collected under informed consent and IRB approval at the University of California, San Francisco. No new data from human or animal subjects were collected.

%% ========== METHODS ==========

\section{Methods}

\subsection{Datasets}
\label{sec:datasets}

Five publicly available datasets were analyzed (Table~\ref{tab:key_resources}). The Steinmetz Neuropixels recordings (DANDI:000017; \citealp{steinmetz2019distributed, zatkahaas2021}) provide spike counts from 12 brain regions during a two-alternative visual discrimination task (licking response) with region-resolved optogenetic silencing ground truth. The ALM photoinhibition data (DANDI:000007, 000009; \citealp{svoboda_alm}) provide paired control and photostimulation sessions in anterior lateral motor cortex during a delayed-response licking task (32 sessions from the primary cohort, 52/46 control/stim sessions from the replication cohort). The human ECoG speech production dataset (DANDI:000019; \citealp{dandi000019}) provides field-potential recordings from 5 cortical areas across 19 sessions in 5 participants during a consonant-vowel-consonant production task. The International Brain Laboratory Brain-Wide Map \citep{ibl2025brainwide} provides Neuropixels recordings from 206 brain regions during a wheel-turning task (86 sessions, 8 regions overlapping with Steinmetz silencing data). The Allen Institute Visual Behavior Neuropixels dataset \citep{siegle2021survey} provides recordings from 49 brain regions during a go/no-go change-detection task (20 sessions).

\begin{table}[t]
\centering
\caption{Key resources.}
\label{tab:key_resources}
\small
\begin{tabular}{p{2.8cm}p{3.2cm}p{2.5cm}p{3.5cm}}
\toprule
\textbf{Resource type} & \textbf{Designation} & \textbf{Source} & \textbf{Identifier} \\
\midrule
Dataset & Steinmetz Neuropixels & DANDI Archive & DANDI:000017 \\
Dataset & ALM photoinhibition & DANDI Archive & DANDI:000007, 000009 \\
Dataset & Human ECoG speech & DANDI Archive & DANDI:000019 \\
Dataset & IBL Brain-Wide Map & IBL data portal & \url{https://internationalbrainlab.org/data} \\
Dataset & Allen VBN & Allen Institute & \url{https://allensdk.readthedocs.io} \\
Software & Analysis code & This paper & \url{https://github.com/elliottower/bracket-norm} \\
Software & Python 3.11 & python.org & RRID:SCR\_008394 \\
\bottomrule
\end{tabular}
\end{table}

\subsection{Silencing effect size}
\label{sec:silencing_def}

The optogenetic silencing effect for each brain region is the fractional decrease in task performance during bilateral photoinactivation, as reported by \citet{zatkahaas2021}: $\Delta_{\text{sil}} = (\text{accuracy}_{\text{control}} - \text{accuracy}_{\text{silenced}}) / \text{accuracy}_{\text{control}}$. Twelve regions have silencing data; 9 of these have sufficient trials per evidence quartile to estimate bracket norm. All correlations with ``silencing importance'' refer to $\Delta_{\text{sil}}$ as the dependent variable.

\subsection{Bracket norm definition}
\label{sec:bn_def}

For each region during the decision window (150--350\,ms post-stimulus), compute the mean choice displacement $\xi = \text{mean}(x_{\text{right}}) - \text{mean}(x_{\text{left}})$ separately within the lowest and highest evidence quartiles ($|\text{contrast}_{\text{L}} - \text{contrast}_{\text{R}}|$), then
\begin{equation}
\text{BN} \;=\; \bigl\lVert \xi(\text{high evidence}) - \xi(\text{low evidence}) \bigr\rVert .
\end{equation}
This approximates the Lie bracket $[\xi_{\text{choice}}, \xi_{\text{evidence}}]$ restricted to the decision window: how much the choice vector field rotates as evidence strength varies (Figure~\ref{fig:quiver}). A region with low bracket norm relays a fixed choice signal regardless of evidence. A region with high bracket norm reshapes its choice representation with evidence---a candidate computation node where the evidence-to-choice transformation takes place. Most individual regions do not significantly exceed a label-shuffled null, but the cross-regional \emph{ranking} predicts causal importance (Appendix~\ref{app:baseline}).

\subsection{Dimensional correction}
\label{sec:dim_correction}

Raw bracket norm is a Euclidean norm in $\mathbb{R}^n$ where $n$ is the number of recorded neurons. After z-scoring, each neuron contributes $O(1)$ variance to the displacement vector $\xi$, so $\lVert \xi \rVert \sim \sqrt{n}$ by the same argument as random-walk step lengths. The scaling is intrinsic to norms in high-dimensional spaces, not a statistical estimation artifact.

We verify this on the Steinmetz Neuropixels dataset \citep{dandi000017}, subsampling $n \in \{10, 20, 30, 50, 75, 100, 150, 200, 250\}$ neurons from the same recording and computing bracket norm at each size. Raw BN scales near-linearly with neuron count ($\rho = +0.989$, $p < 10^{-4}$). The corrected quantity
\begin{equation}
\text{BN}/\!\sqrt{n} \;=\; \frac{\bigl\lVert \xi_{\text{high}} - \xi_{\text{low}} \bigr\rVert}{\sqrt{n}}
\end{equation}
is invariant: mean $0.394 \pm 0.008$ across the full $25\times$ range, coefficient of variation $2.0\%$, correlation with neuron count $\rho = +0.085$ ($p = 0.07$). The rotation angle $\theta = \arccos(\hat{\xi}_{\text{high}} \cdot \hat{\xi}_{\text{low}})$ is also invariant ($\rho = -0.017$), confirming that both dimensionally corrected quantities eliminate the $\sqrt{n}$ artifact (Table~\ref{tab:invariance}).

\begin{table}[t]
\centering
\caption{Synthetic invariance test: correlation of each metric with neuron count across $25\times$ subsampling range \citep{dandi000017}. Raw BN fails; both corrected metrics pass.}
\label{tab:invariance}
\small
\begin{tabular}{lccc}
\toprule
Metric & $\rho$ vs.\ $n$ & $p$ & Verdict \\
\midrule
BN (raw) & $+0.989$ & $<10^{-4}$ & fails \\
BN$/\!\sqrt{n}$ & $+0.085$ & $0.071$ & invariant \\
Rotation angle & $-0.017$ & $0.719$ & invariant \\
\bottomrule
\end{tabular}
\end{table}

$\text{BN}/\!\sqrt{n}$ is an intensive quantity: the per-neuron geometric displacement due to evidence, analogous to temperature (intensive) versus total thermal energy (extensive). Raw BN is the extensive version---it conflates geometric signal with population size.

\subsection{Normalization methods}
\label{sec:normalization}

Two normalization methods are used throughout. Method~B computes $\text{BN}/\!\sqrt{n}$ within each session and then averages across sessions: the unit-of-observation approach. Method~A aggregates raw bracket norm and neuron count across sessions and then normalizes: $\overline{\text{BN}}/\!\sqrt{\bar{n}}$. Because $1/\!\sqrt{n}$ is concave, Jensen's inequality implies Method~A yields a stronger (more negative) correlation with silencing than Method~B. Both methods eliminate the neuron-count confound. We report Method~B throughout for consistency.

\subsection{Statistical methods}
\label{sec:stats}

All correlations are Spearman rank correlations with exact permutation $p$-values (10{,}000 permutations unless stated otherwise). Partial Spearman correlations control for neuron count via rank residuals. Multiple comparisons are corrected by Bonferroni ($\alpha = 0.05 / 20 = 0.0025$ for 20 metrics) rather than false-discovery-rate correction, because the family of tests shares a single null hypothesis (no geometric metric predicts silencing), making family-wise error rate the appropriate quantity. The top-3 matching test uses an exact combinatorial probability: the chance that two independently ranked lists of $n = 9$ items share the same top-3 set is $\binom{9}{3}^{-1} = 1/84 \approx 0.012$. The non-overlapping tail test uses a permutation procedure: BN-ranked top-3 and bottom-3 define two groups; the observed mean silencing difference is compared against 10{,}000 random partitions of the 9 regions into groups of 3 ($p = 0.0008$). Paired comparisons (photoinhibition) use the Wilcoxon signed-rank test. Leave-one-out cross-validation drops each region in turn and recomputes the correlation on the remaining set.

\bibliography{references}

%% ========== SUPPLEMENTARY ==========

\clearpage
\appendix

\section{Robustness and Temporal Specificity}
\label{app:robustness}

\paragraph{Window robustness.} To rule out forking-path dependence on the 150--350\,ms decision window, we recomputed raw bracket norm across eight temporal windows spanning pre-stimulus through post-decision epochs. The silencing correlation is stable: $\rho$ ranges from $+0.41$ (pre-stimulus) to $+0.64$ (mid-stimulus), with no window producing a negative or zero correlation (Table~\ref{tab:windows}). The decision window ($\rho = +0.59$) is neither the strongest nor the weakest.

\begin{table}[h]
\centering
\caption{Bracket norm vs.\ silencing across temporal windows. The correlation is robust to window choice.}
\label{tab:windows}
\small
\begin{tabular}{lcc}
\toprule
Window & $\rho$ vs.\ silencing & $p$ \\
\midrule
Pre-stimulus & $+0.406$ & $0.19$ \\
Early stimulus & $+0.517$ & $0.085$ \\
Mid-stimulus & $+0.636$ & $0.026$ \\
Peri-decision & $+0.608$ & $0.036$ \\
Decision (150--350\,ms) & $+0.587$ & $0.045$ \\
Late decision & $+0.587$ & $0.045$ \\
Post-decision & $+0.608$ & $0.036$ \\
Late & $+0.573$ & $0.051$ \\
\bottomrule
\end{tabular}
\end{table}

\paragraph{Temporal specificity (H27).} Pre-stimulus bracket norm (125--375\,ms, before evidence) already predicts silencing (partial $\rho = +0.533$), and the decision window adds $+0.235$ beyond it (decision-window partial $\rho = +0.768$). The metric captures both a structural baseline (geometric expressiveness present before stimulus) and a decision-specific rotation. ORB contributes decision-specific rotation (pre-stim BN $0.160 \to$ decision $0.188$); ACA is structurally expressive but adds none (pre-stim $0.169 \to$ decision $0.073$).

\paragraph{Leave-one-region-out jackknife.} Dropping each region in turn, the partial $\rho$ ranges from $+0.386$ (dropping PL) to $+0.987$ (mean $+0.668 \pm 0.161$). All 9 jackknife estimates are positive---no single region drives the result.

\section{Joint Predictor Regression}
\label{app:joint}

To assess whether bracket norm carries information beyond neuron count alone, we fit OLS regressions predicting silencing importance from subsets of predictors across the 9 regions with sufficient data.

\begin{table}[h]
\centering
\caption{Variance explained by predictor subsets (OLS, $n = 9$ regions).}
\label{tab:joint}
\begin{tabular}{lc}
\toprule
Predictor(s) & $R^2$ \\
\midrule
Bracket norm alone & 0.652 \\
Neuron count alone & 0.580 \\
Choice probability alone & 0.004 \\
BN + NC + CP + Accuracy & \textbf{0.809} \\
\bottomrule
\end{tabular}
\end{table}

Bracket norm alone explains 65.2\% of the variance in silencing importance, exceeding neuron count (58.0\%). Choice probability contributes nothing ($R^2 = 0.004$). The full model ($R^2 = 0.809$) shows bracket norm and neuron count carry partially overlapping but not redundant information.

\section{Mimic Mechanism Discriminator}
\label{app:mimic}

To distinguish whether bracket norm marks evidence-integration nodes or reflects attractor persistence, we correlated bracket norm with two dynamical signatures across all 36 Steinmetz regions (Table~\ref{tab:mimic}).

\begin{table}[h]
\centering
\caption{Mimic mechanism discriminator: bracket norm is orthogonal to both dynamical signatures, ruling out the attractor-persistence interpretation. Eigenspectrum slope independently predicts silencing importance.}
\label{tab:mimic}
\small
\begin{tabular}{lccc}
\toprule
Comparison & $\rho$ with BN & $p$ & Partial $\rho$ vs.\ silencing \\
\midrule
Autocorrelation timescale $\tau$ & $+0.231$ & $0.17$ & --- \\
Eigenspectrum slope & $-0.037$ & $0.83$ & $-0.776$ \\
\bottomrule
\end{tabular}
\end{table}

Bracket norm is orthogonal to both $\tau$ and eigenspectrum slope. Eigenspectrum slope independently predicts silencing after controlling neuron count (partial $\rho = -0.776$): steeper spectral falloff (more low-rank activity) associates with greater causal importance once recording yield is removed, consistent with the main-text finding that compact, evidence-sensitive coding characterizes causally important regions.

\section{Baseline Separation Test}
\label{app:baseline}

We tested whether bracket norm exceeds chance (label-shuffled null) at each region, running 1{,}000 permutations per region across all 36 Steinmetz regions.

Only 3 of 36 regions show bracket norm significantly above the shuffle null ($p < 0.05$): LS ($p = 0.012$, ratio $= 1.38$), PL ($p = 0.035$, ratio $= 1.32$), and VISpm ($p = 0.028$, ratio $= 1.20$). Across all regions the mean real-to-shuffle ratio is $1.049$---bracket norm is on average only 5\% above chance. This is not a weakness: the cross-regional \emph{ranking} of bracket norm predicts causal importance, even though most individual regions do not significantly exceed baseline. The predictive signal is in the relative ordering, not in absolute departure from chance.

\section{Allen VBN Per-Region Rankings}
\label{app:vbn}

\begin{table}[h]
\centering
\caption{Allen VBN: bracket norm by region (trial-length evidence, 20 sessions). Top 10 and bottom 5 of 42 regions shown.}
\label{tab:vbn_regions}
\small
\begin{tabular}{lcccc}
\toprule
Region & BN mean & BN std & Neuron count & Sessions \\
\midrule
APN & 43.5 & 18.7 & 205 & 17 \\
MRN & 26.8 & 15.0 & 99 & 16 \\
CA1 & 26.3 & 10.0 & 274 & 19 \\
SNr & 25.9 & 2.2 & 30 & 2 \\
MB & 19.5 & 12.2 & 57 & 18 \\
VISp & 18.2 & 6.1 & 163 & 20 \\
VISpm & 18.1 & 13.8 & 125 & 20 \\
DG & 17.5 & 8.8 & 143 & 19 \\
VISrl & 17.3 & 5.9 & 111 & 20 \\
ZI & 16.8 & 6.2 & 55 & 7 \\
\midrule
POL & 10.4 & 5.7 & 47 & 12 \\
POST & 10.2 & 6.5 & 39 & 11 \\
ProS & 9.2 & 4.3 & 41 & 8 \\
DT & 6.5 & 0.0 & 20 & 2 \\
HY & 5.0 & 0.0 & 10 & 1 \\
\bottomrule
\end{tabular}
\end{table}

The highest-bracket-norm regions in Allen VBN are subcortical and midbrain structures (APN, MRN, SNr, MB, ZI) alongside hippocampal (CA1, DG) and primary visual cortex. The ranking is anatomically plausible: anterior pretectal nucleus (APN) and midbrain reticular nucleus (MRN) sit at sensorimotor transformation junctions where choice encoding should rotate most with evidence strength.

\section{IBL Per-Region Detail}
\label{app:ibl}

\begin{table}[h]
\centering
\caption{IBL: per-region detail for the 8 regions overlapping with Steinmetz silencing data.}
\label{tab:ibl_detail}
\small
\begin{tabular}{lccccc}
\toprule
Region & IBL BN & Rotation (rad) & Principal angle ($^\circ$) & Selectivity & Within-region CV \\
\midrule
ACA & 4.36 & 0.98 & 34.5 & 0.044 & 0.47 \\
RSP & 4.21 & 1.41 & 47.4 & 0.041 & 0.48 \\
MOs & 3.62 & 0.99 & 51.5 & 0.041 & 0.27 \\
VISl & 3.51 & 1.31 & 51.5 & 0.033 & 0.68 \\
VISam & 3.51 & 1.37 & 45.8 & 0.037 & 0.46 \\
VISpm & 3.46 & 1.32 & 50.6 & 0.022 & 0.40 \\
VISp & 3.19 & 1.37 & 51.2 & 0.023 & 0.47 \\
PL & 2.92 & 1.15 & 41.6 & 0.042 & 0.26 \\
\bottomrule
\end{tabular}
\end{table}

The key diagnostic is the contrast between columns: rotation angles are large (0.98--1.41~rad, comparable to Steinmetz) and principal angles indicate distinct subspaces (34--52$^\circ$), yet bracket-norm values span only 2.92--4.36 (CV $= 0.15$). The geometry bracket norm measures is present in IBL; the dynamic range to rank regions by it is absent due to weak evidence-quartile separation from discrete contrast levels.

\section{Radial--Angular Decomposition}
\label{app:decomposition}

The bracket vector $\xi_{\text{high}} - \xi_{\text{low}}$ can be decomposed into a radial component (magnitude change along the mean choice axis) and an angular component (rotation orthogonal to it). Let $\hat{m} = (\xi_{\text{high}} + \xi_{\text{low}}) / \lVert \xi_{\text{high}} + \xi_{\text{low}} \rVert$. Then
\[
\text{BN}_{\text{rad}} = |\langle \xi_{\text{high}} - \xi_{\text{low}}, \hat{m} \rangle|, \qquad
\text{BN}_{\text{ang}} = \lVert (\xi_{\text{high}} - \xi_{\text{low}}) - \langle \xi_{\text{high}} - \xi_{\text{low}}, \hat{m} \rangle \hat{m} \rVert .
\]

Across the 12 silencing regions in Steinmetz, the angular component dominates at approximately 4:1 ratio (mean angular/radial $\approx 3.8$), and the mean rotation angle between $\xi_{\text{low}}$ and $\xi_{\text{high}}$ is $85$--$90^\circ$. Evidence does not merely scale the choice vector; it rotates it into a nearly orthogonal direction at high evidence. The angular component after $\sqrt{n}$ correction correlates with silencing at $\rho = +0.587$ ($p = 0.045$), while the radial component shows $\rho = +0.517$ ($p = 0.085$)---both contribute, but the directional change carries more of the signal.

On Svoboda ALM data with photostimulation \citep{dandi000007, dandi000009} (44 sessions with both conditions), the principal angle between control and photostim bracket vectors is $85.1^\circ \pm 17.7^\circ$ (median $84.5^\circ$), indicating photostim rotates the bracket direction into a nearly orthogonal subspace. Within each condition, the rotation angle between low- and high-evidence choice vectors is comparable (ctrl $81.0^\circ$, stim $79.9^\circ$). The angular/radial ratio shifts from $17.0$ (control) to $4.6$ (stim)---photostim increases the radial component disproportionately while both conditions retain angular dominance. This produces a triple dissociation: magnitude (BN increases $+59\%$), within-condition rotation (preserved at ${\sim}80^\circ$), and bracket axis (rotated $85^\circ$). Perturbation changes \emph{which neurons} carry the evidence-dependent signal without disrupting \emph{how} evidence modulates choice encoding within the active subspace.

\section{Holonomy Reframing}
\label{app:holonomy}

Bracket norm admits a differential-geometric interpretation as the curvature of a discrete connection on a choice bundle over evidence space. Partition evidence into $Q = 4$ quartiles and compute $\xi_q$ for each. The discrete connection coefficient between adjacent quartiles is
\[
\Gamma_{q \to q+1} = \arccos\!\bigl(\hat{\xi}_q \cdot \hat{\xi}_{q+1}\bigr),
\]
and the holonomy (total accumulated rotation) is $\Phi = \sum_{q=1}^{Q-1} \Gamma_q$. The per-step curvature is $\kappa = \Phi / (Q - 1)$.

Across the 12 silencing regions, curvature is uniform ($\kappa \approx 28$--$32^\circ$), while BN/\!$\sqrt{n}$ varies by $4\times$. This dissociates magnitude from direction: all regions rotate the choice vector by similar angles per evidence step, but causally important regions produce larger per-neuron displacements at each step. The curvature uniformity is consistent with the radial/angular decomposition---BN variation across regions is primarily a gain (magnitude) effect, not a rotation (curvature) effect.

\section{Causal Baselines}
\label{app:causal_baselines}

Five standard causal inference methods were tested on evidence-binned pseudo-time-series of $\xi$ (6--10 bins per session, 500 permutations each). All return null results (Table~\ref{tab:causal_baselines}).

\begin{table}[h]
\centering
\caption{Causal inference baselines on Steinmetz data. All five methods return null results due to insufficient time-series length (6--10 evidence bins versus the hundreds of points required for Takens embedding or lag-adjusted $F$-tests).}
\label{tab:causal_baselines}
\small
\begin{tabular}{lcc}
\toprule
Method & Result & \% significant \\
\midrule
CCM \citep{sugihara2012detecting} & mean $r \approx -0.3$ & 0\% \\
PCMCI-lite (Granger on PCs) & 0\% forward significance & 0\% \\
Topological CCM & mixed $z$-scores & $< 10$\% \\
Geographical CCM & all $p > 0.4$ vs.\ silencing & 0\% \\
Causalized CCM & mean $r \approx -0.2$ & $< 10$\% \\
\bottomrule
\end{tabular}
\end{table}

The null results reflect a methodological mismatch: these methods detect causal structure in continuous temporal dynamics requiring hundreds of time points, while bracket norm operates on trial-averaged comparisons natural to cognitive neuroscience experimental designs with discrete evidence levels.

\end{document}
